\documentclass[11pt,a4paper]{ctexart}
\input{/Users/shaoyikai/Desktop/考研/22408-notes-skill/assets/note-preamble.tex}
\renewcommand{\NoteFooter}{22408 一页纸笔记 · 408 计算机组成原理}
\linespread{1.12}
\begin{document}
\notetitle{TLB·页表·Cache 变换}{408}{计组·存储系统·虚拟存储与Cache}
\sechead{一、两个地址是同一副骨架}

\[
\text{虚拟地址}=\underbrace{\text{虚页号}}_{\text{查 TLB/页表}}\;\|\;\text{页内偏移}\qquad \text{物理地址}=\underbrace{\text{页框号}}_{\text{查 Cache}}\;\|\;\text{页内偏移}
\]

\warnline{\textbf{偏移位数由页大小定、全程不变}：$4\mathrm{KB}=2^{12}$ 故偏移 12 位，虚页号、页框号都靠它切。}

\sechead{二、一条地址走完三级}

\begin{center}\small
\begin{tabularx}{\linewidth}{@{}>{\raggedright\arraybackslash}X>{\raggedright\arraybackslash}X>{\raggedright\arraybackslash}X>{\raggedright\arraybackslash}X@{}}
\toprule
步骤 & 用什么查什么 & 依据字段 & 未命中的后果 \\
\midrule
① TLB & 虚页号 $\to$ 页框号 & TLB 标记 $+$ 组号 & 查页表，多一次访存 \\
② 页表 & 虚页号 $\to$ 页框号 & 页表基址 $+$ 虚页号 $\times$ 表项长 & 有效位 $0\Rightarrow$ 缺页中断 \\
③ 拼地址 & 页框号 $+$ 页内偏移 & 偏移照抄 & —— \\
④ Cache & 物理地址 $\to$ 数据 & 标记 $+$ 行号 $+$ 块内 & 有效位 $1$ 且标记相等才命中 \\
⑤ 访存 & Cache 未命中才去主存 & —— & 顺带把块调入 Cache \\
\bottomrule
\end{tabularx}
\end{center}

\notebox{TLB 与页表只管「页号 $\to$ 页框号」，偏移一个比特都不动；拼出物理地址后才轮到 Cache——\textbf{Cache 认的是物理地址}。}

\sechead{三、本题的位数（2011 年 408 第 44 题）}

\begin{center}\small
\begin{tabularx}{\linewidth}{@{}>{\raggedright\arraybackslash}X>{\raggedright\arraybackslash}X>{\raggedright\arraybackslash}X>{\raggedright\arraybackslash}X@{}}
\toprule
地址 & 字段（高位 $\to$ 低位） & 位数 & 怎么来的 \\
\midrule
虚址 24 位 & 虚页号 A23\textasciitilde{}A12 & $12$ & $24-12$ \\
 & 页内偏移 A11\textasciitilde{}A0 & $12$ & 页 $4\mathrm{KB}=2^{12}$ \\
物理 20 位 & 页框号 A19\textasciitilde{}A12 & $8$ & $20-12$ \\
 & 页内偏移 A11\textasciitilde{}A0 & $12$ & 同上，不变 \\
Cache & 标记 A19\textasciitilde{}A8 & $12$ & $20-5-3$ \\
 & 行号 A7\textasciitilde{}A5 & $3$ & $8$ 行 $=2^{3}$ \\
 & 块内 A4\textasciitilde{}A0 & $5$ & 块 $32\mathrm{B}=2^{5}$ \\
\bottomrule
\end{tabularx}
\end{center}

\begin{itemize}

\item $16\mathrm{MB}=2^{24}$、$1\mathrm{MB}=2^{20}$；\textbf{标记 $=$ 地址位数 $-$ 行号 $-$ 块内}。自查：题给标记 \texttt{020/01D/105} 是 3 位十六进制 $=12$ 位，与算出吻合即切对。

\end{itemize}

\sechead{四、四问逐条解}

\textbf{(3) 001C60H}：高 12 位 $001\mathrm{H}$ 查页表得「有效位 $1$、页框号 $04\mathrm{H}$」$\Rightarrow$ \textbf{在主存}，物理地址 $=04\mathrm{H}\|C60\mathrm{H}=04C60\mathrm{H}$；拆出标记 $04C\mathrm{H}$、行号 $3$，第 3 行标记 $105\mathrm{H}\ne04C\mathrm{H}$ $\Rightarrow$ \textbf{不命中}。

\textbf{(4) 024BACH}：$8/4=2$ 组 $\Rightarrow$ 组号 1 位、标记 $12-1=11$ 位。虚页号 $024\mathrm{H}=0000\,0010\,0100\mathrm{B}$，\textbf{最低位 $0$ 是组号}，高 11 位 $012\mathrm{H}$ 是标记；组 0 中有「有效位 $1$、标记 $012\mathrm{H}$」的项 $\Rightarrow$ \textbf{TLB 命中，该页在主存}。

\warnline{组号取虚页号\textbf{最低位}：$024\mathrm{H}$ 是偶数页 $\to$ 组 0。这一步错的人最多。}

\sechead{五、易错点}

\warnline{① \textbf{偏移不变}。物理地址不是"重算偏移"，而是页框号替换虚页号、后 12 位照抄。}

\warnline{② \textbf{三种有效位别混}：页表项有效位 $=$ 页在主存；TLB 项有效位 $=$ 这条映射可用；Cache 行有效位 $=$ 这块数据可用。三者互相独立。}

\warnline{③ \textbf{TLB 命中 $\ne$ Cache 命中}，也不等于必然不缺页——只有 TLB 项有效位为 $1$ 才保证页在主存。}

\warnline{④ 直接映射行号取物理地址\textbf{中间 3 位}，标记在高位；组相联组号位数 $=\log_2(\text{项数}/\text{路数})$。}

\begin{itemize}

\item 口诀：\textbf{拆页号、查 TLB、偏移照搬、拼物理、再拆标记行号块内}。

\end{itemize}

\end{document}
